\scanpage{353}
\mspage{8}
above,
\[
\Lambda^{M}(x_{1},\,x_{2},\,\ldots : d)\to A_{M}.
\]

\margin{\handwritten{\(\unsure{\vert|1}\)}}\uline{Theorem} \uline{D}: \quad (Generalized de Rham) The algebraic model of the \(C^{\infty}\) de Rham complex \(\Lambda^{M}(x_{1},\,x_{2},\,\ldots : d)\) is isomorphic to the model of the rational homotopy type \(\Lambda(x_{1},\,x_{2},\,\ldots : d)\) tensored by the real numbers.

This theorem is proved by choosing a smooth triangulation\unsure{'} and considering the diagram of algebras.

\diagram{\begin{tikzcd}[ampersand replacement=\&]
\& \{\text{polynomial forms}\}\otimes R \arrow[d] \\
\{\text{smooth forms}\}\arrow[r] \& \{\text{piecewise smooth forms}\}
\end{tikzcd}}

Each of these algebras computes the cohomology of \(M\) so the three models are isomorphic by uniqueness.

So we can compute the homotopy groups of a simply connected manifold tensored by the reals, the Whitehead products and so on. For non-simply connected manifolds we can determine from the smooth de Rham complex in dimensions 1 and 2 (as described in Appendix N) a tower of real nilpotent Lie groups \(\{\ldots\,N_{k}\to N_{k-1}\to\ldots\to R^{n}\to e\}\) which are related to the fundamental group as follows: if \(\Gamma_{k}\) denotes \(\pi_{1}\)(modulo \(k\)-fold commutators then modulo torsion), then the geometry of \(M\) determines an embedding of \(\Gamma_{k}\) in \(N_{k}\) as a discrete subgroup with compact quotient. \(N_{k}\) is determined up to isomorphism by this property.
% archive sheet number: 354
